\documentclass[11pt]{article}
\usepackage[margin=1in]{geometry}
\usepackage{amsmath,amssymb,mathtools}
\usepackage{graphicx}
\usepackage{booktabs}
\usepackage{microtype}
\usepackage[numbers,sort&compress]{natbib}
\usepackage[colorlinks=true,allcolors=blue]{hyperref}

\newcommand{\Oh}{\mathcal{O}}
\newcommand{\Kgg}{K_{gg}}
\newcommand{\Kgp}{K_{gp}}
\newcommand{\Gfour}{\mathcal{G}^{\mathrm{eff}}_4}
\newcommand{\PL}{P_L}
\newcommand{\PR}{P_R}

\title{Fourth-Order Restoration of the Canonical Gauge Sector\\in Refined Four-Dimensional Regge Calculus}
\author{Craig Botnen}
\date{}

\begin{document}
\maketitle

\begin{abstract}
Discretizations of general relativity generically break diffeomorphism symmetry, replacing exact canonical constraints by pseudo-constraints away from flat or specially curved sectors. Earlier Regge-calculus work established exact and approximate Bianchi identities, contracted conservation laws in the small-deficit regime, and quadratic lifting of would-be gauge Hessian modes with deficit angle. We study how this mechanism propagates through canonical tent-move evolution under shape-regular refinement of a four-dimensional Freudenthal--Kuhn (FK) length-Regge family.

For stationary near-flat Euclidean Regge data, a contracted-Bianchi argument implies that the corrected vertex-displacement sector of the reduced middle-slice Hessian has no contributions at orders $h^2$ or $h^3$. After eliminating physical mixing, the leading gauge obstruction is the fourth-order Schur term
\[
\Gfour=\PL\bigl(H_4-H_2H_0^+H_2\bigr)\PR.
\]
Under a uniform physical spectral gap and fourth-order transversality, this yields an $\Oh(h^4)$ gauge-sector pseudo-constraint defect. On fully stationary ordinary-Regge FK two-step complexes containing 48 four-simplices, the four lifted gauge singular values of the middle-slice Hessian scale as $h^4$ while the 12-dimensional physical gap remains $\Oh(1)$. Differentiating the stationary equations to form the old/new mixed Lagrangian Hessian $K$, we find the block hierarchy
\[
\Kgg=\Oh(h^4),\qquad \Kgp=\Oh(h^2).
\]
The stationary response is bounded on boundary deformation directions but grows approximately as $h^{-2}$ for unrestricted physical boundary forcing.

Independently, direct normal-deformation commutators on neighboring FK tetrahedra recover the inverse-metric structure function
\[
\beta^a=\sigma g^{ab}(N\partial_bM-M\partial_bN),
\]
with $g^{-1}$ fixed by the FK graph symbol. A fixed-volume weak assembly preserves the fourth-order deformation-sector rate. The result is an on-shell, perturbative Euclidean length-Regge statement; it does not assert exact finite-spacing diffeomorphism symmetry or an unrestricted off-shell Dirac algebra.
\end{abstract}

\section{Introduction}
The hypersurface-deformation algebra (HDA) is the canonical imprint of spacetime covariance in general relativity. In continuum ADM variables, the bracket of two normal deformations closes into a tangential deformation with inverse spatial metric as a structure function. Generic discretizations do not preserve this symmetry exactly. In Regge calculus, flat backgrounds possess vertex-displacement symmetries, whereas curvature generally lifts those directions and produces pseudo-constraints \cite{BahrDittrich2009,DittrichHoehn2010}.

The refinement mechanism has several known ingredients. Hamber and Kagel derived an exact nonlinear Regge Bianchi identity from products of finite hinge rotations \cite{HamberKagel2004}. Gentle \emph{et al.} formulated a Kirchhoff-like contracted conservation law and power-counted its refinement behavior \cite{GentleEtAl2009}. Williams related the linearized contracted Bianchi identities directly to combinations of the four-dimensional Regge equations \cite{Williams2012}. Bahr and Dittrich further observed that small Hessian eigenvalues associated with broken vertex-displacement symmetry vanish quadratically with the deficit angle \cite{BahrDittrich2009}. Since shape-regular refinement of a smooth metric gives deficits of order $h^2$, these results anticipate a fourth-order gauge-lifting scale.

What is not supplied by those observations alone is the corresponding statement after physical/gauge reduction of the canonical old/new mixed Lagrangian Hessian. The full canonical map need not be fourth-order close to a symmetry: physical boundary perturbations can generate parametrically larger responses. Our contribution is therefore a block statement. On fully stationary FK two-step length-Regge configurations we obtain
\[
H_{gg}=\Oh(h^4),\qquad \Kgg=\Oh(h^4),\qquad \Kgp=\Oh(h^2),
\]
with a bounded stationary response only on the geometric boundary-deformation subspace.

This on-shell formulation is deliberate. Modern analyses emphasize that geometric hypersurface deformations and canonical Poisson-bracket generators are not a purely kinematical off-shell equivalence \cite{BojowaldDuqueShah2025}. Our theorem is therefore stated for stationary middle-slice Regge data, followed by a canonical corollary restricted to deformation sources.

\section{FK metric dictionary and deformation bracket}
The seven positive FK spatial directions are
\[
(1,0,0),\ (0,1,0),\ (0,0,1),\ (1,1,0),\ (1,0,1),\ (0,1,1),\ (1,1,1).
\]
For equal weights the principal tensor is
\[
A=\begin{pmatrix}4&2&2\\2&4&2\\2&2&4\end{pmatrix},
\qquad g=A^{-1}.
\]
Thus the inverse metric appearing in the HDA is $g^{-1}=A$. The direct neighboring-tetrahedron normal-deformation calculation yields
\[
\beta^i=\sigma g^{ij}(N\partial_jM-M\partial_jN),
\]
where $\sigma=n\cdot n$ is the normal-sign convention. The Euclidean FK controls verify $\sigma=+1$ exactly in the continuum deformation limit. A separate Minkowski-normal calculation verifies the expected sign reversal $\sigma=-1$ kinematically; the fully stationary Lorentzian canonical calculation is not part of the present result.

\section{Stationary middle-slice estimate}
We work with normalized lengths $\hat l=l/h$ and dimensionless action $\widehat S=h^{-2}S$. Let
\[
\widetilde S_h(l^{n-1},l^n,l^{n+1})
\]
be the two-step effective length-Regge action after legitimate local stationary reductions. The middle-slice variables satisfy
\[
E_h:=\frac{\partial\widetilde S_h}{\partial l^n}=0,
\qquad
H_h:=\frac{\partial^2\widetilde S_h}{\partial l^n\partial l^n}.
\]
For a normalized vertex-displacement field $Y_A$, define $B_A=Y_A\cdot E_h$. Then
\[
Y_B[B_A]=(\nabla_{Y_B}Y_A)\cdot E_h+Y_A^TH_hY_B.
\]
At middle-slice stationarity the first term vanishes.

For a hinge triangle $h=(v,a,b)$, let $u=a-v$, $w=b-v$, $d_h=a-b$, and $U_h=(u\wedge w)/|u\wedge w|$. The exact area-gradient identity is
\[
\nabla_vA_h=\frac12 U_h\cdot d_h.
\]
Using the Schlaefli-reduced Regge equations, the vertex contraction is therefore a deficit-bivector/edge contraction. The linear term is removed by the contracted Regge Bianchi identity, while the exact identity contains quadratic and higher curvature contributions \cite{HamberKagel2004,GentleEtAl2009,Williams2012}. Hence $B_A=\Oh(\epsilon^2)$. With smooth-refinement scaling $\epsilon=\Oh(h^2)$ and $\delta_Y\epsilon=\Oh(h^2)$,
\[
Y_A^TH_hY_B=\Oh(h^4).
\]
Writing
\[
H_h=H_0+h^2H_2+h^3H_3+h^4H_4+\Oh(h^5)
\]
therefore gives
\[
\PL H_2\PR=0,\qquad \PL H_3\PR=0.
\]

\section{Schur-corrected fourth-order obstruction}
In a physical/gauge adapted basis,
\[
H_0=\begin{pmatrix}A_0&0\\0&0\end{pmatrix},
\qquad \sigma_{\min}(A_0)\ge\gamma>0.
\]
The corrected gauge vectors satisfy
\[
X_R^{(2)}=-H_0^+H_2X_R^{(0)},\qquad
X_R^{(3)}=-H_0^+H_3X_R^{(0)},
\]
with analogous left corrections. The leading gauge obstruction is
\[
\boxed{\Gfour=\PL(H_4-H_2H_0^+H_2)\PR.}
\]
The $h^3$ curvature-gradient term modifies the corrected vectors and higher remainder but does not enter $\Gfour$. A uniform physical gap implies a bounded pseudoinverse, so the Schur backreaction cannot generate inverse powers of $h$.

\paragraph{Failure map.}
The assumptions are rate-critical rather than cosmetic. If the middle-slice residual is only $\|E_h\|=\Oh(h^q)$, then generically the gauge-row estimate becomes $\Oh(h^{\min(4,q)})$. If the physical gap closes as $\gamma_h\sim h^\alpha$, the Schur term can degrade to $\Oh(h^{4-\alpha})$. Finally, if a boundary source enters the reduced gauge equation at $\Oh(h^r)$ while the gauge Jacobian is $\Oh(h^4)$, the gauge response scales as $\Oh(h^{r-4})$.

\section{Canonical gauge-sector bridge}
Differentiating the stationary middle-slice equations with respect to upper-boundary data gives the standard tent-move identity \cite{BahrDittrich2009}
\[
H_h\mathscr L_h=-K_h,
\]
where $K_h$ is the old/new mixed Lagrangian Hessian and $\mathscr L_h$ is the stationary response. Let $Q_g$ span upper geometric vertex-deformation directions. The reduced gauge equation has the schematic hierarchy
\[
\Psi_h(g,b_g,b_p)=h^4\Psi_4(g,b_g)+h^2S_2b_p+\Oh(h^5).
\]
If $D_g\Psi_4=\Gfour$ is nonsingular, the rescaled implicit-function theorem gives
\[
\|\mathscr L_hQ_g\|=\Oh(1).
\]
In contrast, a physical source entering at $h^2$ may produce $\Oh(h^{-2})$ response. Therefore
\[
X_L^TK_hQ_g=-X_L^TH_h\mathscr L_hQ_g=\Oh(h^4),
\]
which is the deformation-sector canonical statement. No fourth-order claim is made for $X_L^TK_hQ_p$.

\section{Fully stationary FK48 validation}
We solve the ordinary length-Regge equations on a two-step FK complex containing 48 four-simplices. At $h=0.09,0.065,0.045$, all 16 middle/internal equations are solved to approximately $10^{-15}$.

The middle Hessian exhibits a stable 12-dimensional physical sector with smallest singular value approaching $3.982$, while the four lifted gauge singular values fit exponents
\[
3.99931,\quad 3.99933,\quad 4.00188,\quad 4.00651.
\]
The reduced fourth-order gauge Jacobian is full rank, with singular values divided by $h^4$ approaching approximately
\[
0.1696,\quad0.1389,\quad0.1068,\quad0.05368.
\]

For the canonical mixed matrix, the refined finite-difference calculation gives
\[
\|\Kgg\|\propto h^{4.00028},\qquad
\|\Kgp\|\propto h^{1.99764}.
\]
The four singular-value exponents of $\Kgg$ are approximately
\[
4.00028,\quad3.99826,\quad3.99923,\quad4.01677.
\]
The deformation-restricted response remains near $0.333$, while the unrestricted response scales as approximately $h^{-2.0065}$.

\begin{figure}[t]
\centering
\includegraphics[width=0.72\linewidth]{figures/figure1_canonical_block_hierarchy.pdf}
\caption{Direct canonical block hierarchy on the fully stationary FK48 family. The deformation-sector block $K_{gg}$ is fourth order, while gauge-to-physical mixing remains second order.}
\label{fig:blockhierarchy}
\end{figure}

\begin{figure}[t]
\centering
\includegraphics[width=0.72\linewidth]{figures/figure2_Kgg_singular_values.pdf}
\caption{All four singular values of the old/new gauge-to-gauge mixed block $K_{gg}$.}
\label{fig:kggsv}
\end{figure}

\begin{figure}[t]
\centering
\includegraphics[width=0.72\linewidth]{figures/figure3_stationary_response.pdf}
\caption{Stationary-response hierarchy. The response to geometric deformation sources is bounded, while unrestricted physical forcing produces an approximately $h^{-2}$ response.}
\label{fig:response}
\end{figure}

\section{Weak periodic assembly}
For bounded $C^1$ smearings and a stable interpolation, define the cell-integrated defect
\[
\mathcal E_h[N,M]=h^3\sum_{v\in\Lambda_h}e_v[N,M].
\]
Shape regularity gives $\Oh(h^{-3})$ cells in a fixed physical volume. Hence a local $\Oh(h^4)$ deformation-sector coefficient remains $\Oh(h^4)$ after fixed-volume weak assembly:
\[
|\mathcal E_h[N,M]|\le C V\|N\|_{C^1}\|M\|_{C^1}h^4+\Oh(h^5).
\]
This is a statement in this weak, cell-integrated norm only. Rigorous Regge-curvature convergence results motivate the use of shape-regular refinement hypotheses \cite{Christiansen2024,GawlikNeunteufel2025}.

\section{Relation to prior work and limitations}
Quadratic deficit-angle lifting of would-be gauge modes is prior art \cite{BahrDittrich2009}; the present contribution is the propagation of that mechanism through the stationary physical/gauge Schur reduction into the direct canonical gauge-to-gauge block, together with the FK structure-function and weak-assembly checks. Exact discrete HDA constructions in four dimensions are known in flat or homogeneously curved sectors \cite{BonzomDittrich2013}; our result instead quantifies refinement restoration in a generic near-flat curved length-Regge family.

The result does not extend automatically to area Regge calculus, where non-metric modes alter the broken-diffeomorphism structure. It also does not use raw residual errors obtained by simply sampling a continuum solution on a lattice; the principal FK48 validation is performed on fully stationary discrete Regge solutions.

The Lorentzian normal-deformation kinematics has been checked separately and gives the expected $\sigma=-1$ sign. A fully stationary Lorentzian canonical FK48 calculation requires a specified causal sector and is left for future work.

\section{Conclusion}
The FK refinement family exhibits fourth-order restoration specifically in the canonical deformation/gauge sector. On fully stationary ordinary length-Regge solutions,
\[
H_{gg}=\Oh(h^4),\qquad \Kgg=\Oh(h^4),\qquad \Kgp=\Oh(h^2),
\]
with bounded deformation-restricted response and the correct inverse-metric structure function. The block separation is essential: refinement restores the canonical gauge sector without forcing the entire discrete canonical map to become fourth-order close to a gauge symmetry.

\bibliographystyle{unsrtnat}
\bibliography{references}
\end{document}
